% Generated by tools/build_critical_review.py; edit overview.tex for the brief.
\documentclass[11pt]{article}
\usepackage[a4paper,margin=25mm]{geometry}
\usepackage{amsmath,amssymb,amsthm,mathtools,booktabs,microtype,array,xurl}
\setlength{\emergencystretch}{2em}
\usepackage[hidelinks]{hyperref}
\newtheorem{theorem}{Theorem}
\newtheorem{lemma}[theorem]{Lemma}
\newtheorem{corollary}[theorem]{Corollary}
\newtheorem{proposition}[theorem]{Proposition}
\newcommand{\E}{\mathbb E}
\newcommand{\Pp}{\mathbb P}
\newcommand{\Var}{\operatorname{Var}}
\newcommand{\ind}{\mathbf1}
\newcommand{\KL}{\operatorname{KL}}
\newcommand{\Aff}{\operatorname{Aff}}
\newcommand{\Hsq}{\mathsf H^2}
\newcommand{\Xrho}{\mathcal X_\rho}
\title{An iterated-log phase transition in rare-event threshold estimation\\[4pt]\large Specialist review packet}
\author{Project lead: David Iach\\\small AI-assisted research draft}
\date{27 October 2026}
\begin{document}
\maketitle
\section*{Review brief: a possible iterated-log phase transition}
\addcontentsline{toc}{section}{Review brief}

\noindent\textbf{Purpose.} This packet asks for a correctness and prior-art
assessment of a specific statistical claim. The first three pages give the
claim, its vulnerable dependencies, and the verification boundary. The
remaining pages contain the complete existing paper argument with its three
proof inputs expanded in place. This is an AI-assisted research draft; no
external specialist review is claimed.

\subsection*{The experiment and claimed rate}
A fresh coin has hidden magnitude $D$ with density $32d$ on $(0,1/4]$ and
hidden sign $S\mid D=d\sim\operatorname{Ber}(q(d))$. Its flips have mean
$1/2+(2S-1)D$. The unknown measurable allocation satisfies
$0\le q\le1$ and $\theta(q)=\int q\,d\mu\le\rho$, where $0<\rho\le1/4$.
Algorithms know $\mu,\rho,B$, but observe only their independent randomness
and coin flips. They may create, revisit and interleave coins, return biased
nonlinear estimates, and stop at an unbounded time. Every flip is charged.
The budget is $\sup_q\E_qT\le B$, $B\ge1$.

For $L=1+\log(1/\rho)$ and $H=1+\log L$, the paper claims
\[
 R_{\mu,\rho}(B)\asymp
 \min\left\{\rho,\sqrt{\frac{\rho LH}{B}}\right\},
\]
with universal constants for this one fixed magnitude law. The supremum is
over the full global-rarity class; it is not a pointwise constraint
$q(d)\le\rho$, and the target is not the known number $\rho$.

\subsection*{Why this claim matters}
The proposed general variational functional has order
\[
 \Psi_{\mu,\rho}(B^{-1/2})\asymp
 \min\left\{\rho,\sqrt{\frac{\rho L}{B}}\right\}
\]
at this law. If the new lower bound is correct, no universal matching
$C\Psi$ upper bound can hold. For example, the choice
$\rho_n=\exp(-\exp n)$, $B_n=\rho_n^{-2}$ gives an unbounded ratio of
order at least $\sqrt{H_n}$. The optimizer of the variational functional
remains valid; its proposed statistical interpretation is what fails.

The claimed additional cost is the need to discover a coin's sign across
unknown gap scales. Sign-entropy mechanisms are established prior work.
The proposed contribution is their cost-sensitive transfer to the full
population-estimation experiment. No novelty or practical-performance claim
follows merely from this explanation.

\clearpage
\section*{Proof route and questions for the reviewer}
\addcontentsline{toc}{section}{Proof route and review questions}

The lower proof uses the convex subclass
$\mathcal X_\rho=\{0\le q\le1/2,\ \int q\,d\mu\le\rho\}$ and
$J=\rho LH$. Its logical chain is as follows.

\begin{enumerate}
\item \textbf{Null-cost entropy obstruction.} For every one-coin estimator
with uniform bias at most $\rho/64$ and separately finite uniform variance
and work, the claim is
\[
 \left(\sup_{q\in\mathcal X_\rho}\Var_qY\right)\E_0T\ge c_bJ.
\]
The work is measured under the all-negative population. The proof selects
good gaps in separated annuli and restarts the estimator on unused blocks
of the \emph{same} coin to construct one sign test. Finite-prefix
change-of-measure inequalities and overlapping stopping windows supply the
entropy penalty. Fair-coin termination is not assumed.

\item \textbf{One common pair.} A finite categorical reduction and convex
RMS comparison first give a hard pair for each capped local policy.
Revealing the index of a finite policy mixture makes both its null work
and squared Hellinger distance exactly affine in the mixture weights.
Minimax and weak-star compactness are then used to obtain one separated
pair satisfying, for \emph{every} capped policy $A$,
\[
 \mathsf H^2(P^A_-,P^A_+)\le C\delta^2\E_0T_A/J.
\]
The pair is existential. It is not identified with the radial variational
optimizer or selected after observing a transcript.

\item \textbf{Population transfer and termination.} Finite tree leaves
factor by coin even under interleaving. One-coin hybrid simulations use
midpoint tapes for the other coins. The midpoint law dominates half the
null law, so the one-coin bounds can be charged to actual midpoint work.
The same endpoint pair applies at all caps. Finite expected endpoint work
then passes the testing bound to the terminal experiment.
\end{enumerate}

\noindent\textbf{Requested scrutiny.} Does good-gap selection and the
restart construction establish the stated \emph{null-cost} bound? Are the
sufficiency, convexity and compactness hypotheses satisfied strongly enough
to obtain a pair common to all policies? Does the hybrid simulation charge
work under the correct law, including arbitrary randomization and creation?
Does the endpoint limit preserve the full unrestricted model?

Separately, is an equivalent critical population-estimation result already
known? The complete proof credits Chen--Li--Qiao, Juditsky--Nemirovski,
Polyanskiy--Wu and Lee--Valiant. Their ingredients and this particular
application should be assessed separately.

\clearpage
\section*{What has been checked, and what remains open}
\addcontentsline{toc}{section}{Verification and provenance}

\begin{center}
\begin{tabular}{>{\raggedright\arraybackslash}p{0.34\textwidth}>{\raggedright\arraybackslash}p{0.58\textwidth}}
\toprule
Claim or component & Verification boundary\\
\midrule
Original arbitrary-law minimax lower bound & Lean-checked in Tier 6,
with a constructed infinite experiment and uniform expected budget.
It is a different theorem from the critical rare-event result.\\[5pt]
Finite-gap sign entropy & The statistical BC2 theorem is Lean-checked for
actual seeded Bernoulli tests with observable stopping and actual flip
costs. Its application through restarts and good-gap selection remains
unformalized.\\[5pt]
Common-point quantifier interchange & A compact-convex theorem for all
finite mixtures is Lean-checked. The statistical population set and its
mixture-information bounds remain hypotheses to be established.\\[5pt]
Critical minimax theorem & Paper proof with agent review. Its full
statistical dependency chain is not Lean-checked and has no external
specialist review.\\[5pt]
Numerical and exact finite checks & Sanity and reproducibility evidence.
They do not certify the continuum theorem or distinguish its asymptotic
rate from the proposed Psi rate.\\
\bottomrule
\end{tabular}
\end{center}

The existing critical proof was reviewed by an adversarial agent. A second
review disclosed that its author also wrote the base-cost component; it
must not be counted as a fully independent review of that component.
Model diversity is not established. The follow-up audit has not identified
a blocking error in its initial pass, but that observation is not a
certificate of correctness.

\subsection*{Source and reproduction}
The complete mathematical argument below is expanded from the four
source files at published commit
\nolinkurl{b3f88e6bc02b3d59ec9fe99529688e2bd911e1db}.
The accompanying build script verifies their existing SHA-256 lock before
assembly. Editorial changes in this packet concern the title, this brief,
and layout. The inherited mathematical statements and arguments are preserved;
monospaced source paths and hashes are made line-breakable for typesetting.

\begin{itemize}
\item \href{https://github.com/davidiach/adaptive-threshold-estimation/blob/accdb144972edd255dedec50cc1cbcaa45587040/docs/CRITICAL_AUDIT.md}{Follow-up dependency audit and formalization boundary}.
\item \href{https://github.com/davidiach/adaptive-threshold-estimation/tree/b3f88e6bc02b3d59ec9fe99529688e2bd911e1db/research/2026-10-27-nonlinear-critical}{Original proof bundle, source lock, and scoped reviews}.
\item \href{https://github.com/davidiach/adaptive-threshold-estimation/blob/1c1b761702126aa543a78d5e3f5ba23c8a5d87b8/docs/process/METHOD.md}{Process note and limits of the evidence}.
\end{itemize}

\noindent A matching characterization for arbitrary magnitude laws remains
open. This packet concerns a fixed quadratic law. It does not establish
uniformity as a power exponent approaches two, explicit useful constants,
or practical superiority of an implementation.

\clearpage
\part*{Complete critical-law argument}
\addcontentsline{toc}{part}{Complete critical-law argument}


\noindent\textbf{Status.} Complete paper proof with an independent adversarial
agent review and additional component cross-checks. This is not a Lean proof,
an external specialist review, or a novelty certification. Exact reviewed
source identities and findings are maintained separately. The strongest claim below concerns
the full iid fresh-coin experiment, including nonlinear outputs and interaction
among coins; it is not inferred merely from the earlier additive-class theorem.

\section{Experiment, quantifiers, and result}
Throughout fix $a=1/4$ and
\[
 \mu(dd)=\frac{2d}{a^2}\ind_{(0,a]}(d)\,dd.
\]
For $0<\rho\le1/4$, let
\[
 \mathcal Q_\rho=\{q:(0,a]\to[0,1]\text{ measurable}:\int q\,d\mu\le\rho\},
 \qquad \theta(q)=\int q\,d\mu.
\]
Each genuine fresh coin has an independent hidden magnitude $D\sim\mu$,
hidden sign $S\mid D\sim\operatorname{Ber}(q(D))$, and conditionally iid
flips of mean $1/2+(2S-1)D$. Neither latent variable is observed. One realization
uses independent uniforms on $(0,1)$ for each coin: $D=a\sqrt{U_0}$,
$S=\ind_{\{U_1\le q(D)\}}$, and flips
$X_k=\ind_{\{U_{k+1}\le1/2+(2S-1)D\}}$ for $k\ge1$. Measurable functions
equal $\mu$-almost everywhere define the same experiment; Borel versions may
be used in this standard probability space.

The algorithm has an independent seed and sees only its seed and the observed
history. At each step it stops and returns a real measurable output, revisits
an observed coin, or observes a fresh coin. Every flip costs one. The output
may be biased and nonlinear, and the choices may depend on all previous
observations. Stopping may be unbounded. A design may know $\mu,\rho,B$ but is
fixed before the unknown $q$ is supplied. Work $T=\infty$ on nontermination.
Define
\[
 R_{\mu,\rho}(B)=\inf_{A:\,\sup_{q\in\mathcal Q_\rho}\E_qT_A\le B}
       \sup_{q\in\mathcal Q_\rho}\bigl(\E_q(\widehat\theta_A-\theta(q))^2\bigr)^{1/2}.
\]
This class permits arbitrary creation and interleaving, not a globally fixed
number of coins. We number coins by their first charged observation. This
is equivalent to giving names to iid unseen coins earlier: an unused coin's
latent state is independent of the observed history and seed, and choosing
among unused names using only that information still selects an iid fresh
state. The relabeling retains the map for every later revisit. Creating or
discarding an unobserved coin supplies no information. Consequently at most
$K$ different coins occur before the $K$th charged observation.

Write $L=1+\log(1/\rho)$, $H=1+\log L$, and $J=\rho LH$.
For the lower proof we use the convex subclass
\[
 \Xrho=\{q\in L^\infty(\mu):0\le q\le1/2\ \mu\text{-a.e.},\ \int q\,d\mu\le\rho\}.
\]
This pointwise cap restricts the lower-bound adversary only; the algorithmic
upper and minimax definition still use the full global-rarity class.

\begin{theorem}[Full critical law]\label{thm:critical-minimax}
There exist universal constants $0<c<C<\infty$ such that, for the one fixed
law above, every $0<\rho\le1/4$ and $B\ge1$ satisfy
\[
 c\min\{\rho,\sqrt{J/B}\}\le R_{\mu,\rho}(B)
       \le C\min\{\rho,\sqrt{J/B}\}.                    \tag{A}
\]
The lower bound holds for every observable algorithm in the stated full
experiment, including arbitrary randomization, nonlinear estimation and
data-dependent creation, revisiting and interleaving. The upper permits the
known rarity and law; it has separately finite uniform work and squared error.
\end{theorem}

The upper construction was already available in the separately owned critical
note, now at commit \nolinkurl{a8a8dc49115aabdd1fba8d53d3390eb0a5eee864}; it is
reproduced with attribution in the appendix. The new lower proof has three
steps. First, sign-testing entropy forces a small-bias variance product using
\emph{all-negative work}, not a supremum of work over populations. Second,
convex estimation and minimax supply one hard pair valid for every capped
single-coin observation policy. Third, a midpoint simulation charges those
bounds to the actual work of an interactive population algorithm. The last
step is the established Lee--Valiant Hellinger decomposition mechanism
\cite[Lemma 17]{lv}, proved below at finite prefixes to make its observation
and stopping scope explicit.

The convex-estimation step follows Juditsky--Nemirovski's saddle-point
framework \cite{jn}; the passage from its exponential tails to RMS risk is
written out here. Polyanskiy--Wu \cite{pw} already establish broad quadratic-risk
affine comparisons and convex duality for functional estimation. Neither that
general principle nor the Hellinger decomposition is claimed as new. The
cost-sensitive common pair and its application to this critical hidden-coin
class are the proposed contribution; priority remains unverified.

% BEGIN LOCKED SOURCE: base_cost.tex
% Input fragment. Requires amsmath, amssymb, amsthm and the parent theorem environments.
% The mathematical adaptation was constructed by the alternate-analysis agent;
% it is not an independent review of its own contribution.
\section{A single-coin entropy bound with all-negative reference cost}

Fix $a=1/4$ and $\mu(dd)=2d a^{-2}\mathbf1_{(0,a]}(d)\,dd$.
Throughout this section,
\[
 \mathcal X_\rho=\left\{q:0\le q\le\tfrac12,
                    \ \int q\,d\mu\le\rho\right\},
 \qquad 0<\rho\le\tfrac14,
 \qquad L=1+\log(1/\rho),\quad H=1+\log L.
\]
All functions defining populations are measurable. A single genuine coin
has hidden $D\sim\mu$ and $S\mid D=d\sim\operatorname{Ber}(q(d))$;
conditionally on $(D,S)$ its observations are iid Bernoulli with parameter
$1/2+(2S-1)D$. A policy receives only its independent seed and observed
bits. It may know $\mu,a,\rho$, but not $q,D,S$.

\begin{theorem}[The critical product with null-population cost]
\label{thm:base-cost}
There is a universal $c_0>0$ such that the following holds for every
$0<\rho\le1/4$. Let an observable single-coin procedure have real output
$Y$ and stopping time $T\in\mathbb N_0\cup\{\infty\}$, and suppose
\[
 \sup_{q\in\mathcal X_\rho}|\mathbb E_qY-\theta(q)|\le\rho/64,
 \qquad
 V:=\sup_{q\in\mathcal X_\rho}\operatorname{Var}_qY<\infty,
 \qquad
 \sup_{q\in\mathcal X_\rho}\mathbb E_qT<\infty,
 \qquad \theta(q)=\int q\,d\mu.
\]
Writing $W_0=\mathbb E_{q=0}T$, one has
\[
                 V W_0\ge c_0\rho L H.                 \tag{BC1}
\]
The same order lower bound holds with
$M=\sup_{q\in\mathcal X_\rho}\mathbb E_qY^2$ in place of $V$.
The cost is the expectation under the all-negative population, rather
than the supremum of expected work over the parameter class. The policy
may stop before observing the coin and may have unbounded stopping.
\end{theorem}

The proof adapts the separately documented critical-two single-coin
argument. We reproduce its finite-gap dependency to make this statement
self-contained. Its entropy mechanism follows Chen--Li--Qiao,
arXiv:1608.06031, Lemma~4.1; the Bernoulli, one-sided-cost statement used
here is proved explicitly below. No novelty is claimed for that mechanism.

\begin{lemma}[Finite-gap sign entropy]\label{lem:finite-sign-entropy}
There are universal $c_s>0$ and $K_s<\infty$ with the following property.
Let $d_1>\cdots>d_K>0$ satisfy $d_j/d_{j+1}\ge2$ and
$d_j\le1/128$. Suppose one observable sign-testing procedure, the same
at all these gaps, has error at most $1/10$ under every Bernoulli law
$1/2\pm d_j$, and finite expected stopping time $\tau$ under both signs.
For $K\ge K_s$,
\[
              \sum_{j=1}^K d_j^2\mathbb E_{-d_j}\tau
                   \ge c_s K\log K.                    \tag{BC2}
\]
Correctness and termination under the fair coin are not required.
One may take $c_s=1/512$ and $K_s=\lceil e^{32}\rceil$.
\end{lemma}
\begin{proof}
Put $b_0=1/1024$, $n_j=\lfloor b_0/d_j^2\rfloor$,
$u_j=\lceil4\mathbb E_{-d_j}\tau\rceil$,
$c_j=d_j^2\mathbb E_{-d_j}\tau$, and
$\bar c=K^{-1}\sum_jc_j$. Since $d_j\le1/128$,
$n_j\ge b_0/(2d_j^2)\ge1$. The event of returning minus by time $n_j$
is measurable in the independent seed and the first $n_j$ flips, including
unused coordinates if the test has stopped. The KL divergence between
these full prefixes at the minus and plus laws is at most
\[
 n_j\frac{(2d_j)^2}{(1/2+d_j)(1/2-d_j)}
 \le\frac{64}{3}n_jd_j^2\le\frac1{48}.
\]
Bernoulli KL is at most its chi-square divergence. Pinsker's inequality
bounds the total variation by less than $1/8$. Since an early minus is
an error under the plus law, it has probability at most $1/10+1/8$ under
the minus law. Markov's inequality gives
$\mathbb P_{-d_j}\{\tau>u_j\}\le1/4$. Consequently
\[
 A_j=\{\text{output minus and }n_j<\tau\le u_j\}
 \quad\text{satisfies}\quad
 \mathbb P_{-d_j}(A_j)
 \ge1-\tfrac1{10}-(\tfrac1{10}+\tfrac18)-\tfrac14>\tfrac14.
\]

Let $P_0$ denote the fair-coin law with the same seed and policy. Comparing
the deterministic prefixes of length $u_j$, their KL divergence from
minus to fair is at most
$4d_j^2u_j\le16c_j+1$. For an event of probability $p\ge1/4$ under
the minus law and probability $z$ under the fair law, binary data
processing and $\operatorname{kl}(p,z)\ge p\log(1/z)-\log2$ therefore give
\[
                     P_0(A_j)\ge e^{-64c_j-8}.          \tag{BC3}
\]
Only deterministic finite prefixes have been compared; this does not
invoke optional stopping under a possibly nonterminating fair law.

At least $K/2$ indices form a set $G$ with $c_j\le2\bar c$. For them,
\[
 (n_j,u_j]\subset
 [b_0/(2d_j^2),(8\bar c+1)/d_j^2].
\]
The scales $d_j^{-2}$ have successive ratios at least four. At any finite
time the number of these intervals containing that time is at most
$1+\log_4((16\bar c+2)/b_0)$. All events $A_j$ refer to the same terminal
time and output on a realization under the fair law; on a nonterminating
path none occurs. Summing their indicators and using (BC3) yields
\[
 (K/2)e^{-128\bar c-8}
        \le1+\log_4(1024(16\bar c+2)).                  \tag{BC4}
\]
If $x=\log K\ge32$ and $\bar c<x/512$, the left side exceeds
$\tfrac12e^{3x/4-8}\ge\tfrac12e^{x/2}\ge x^2/16$.
The right side is less than
$1+\log_4(32x+2048)<6+x/4$, contradicting $x\ge32$.
Thus $\bar c\ge(\log K)/512$, as claimed.
\end{proof}

\begin{proof}[Proof of Theorem~\ref{thm:base-cost}]
The bias bound implies $M\le V+(65\rho/64)^2<\infty$.
First define the nonnegative extended conditional moments
\[
 M_s(d)=\mathbb E_{d,s}Y^2,\qquad w_s(d)=\mathbb E_{d,s}T.
\]
The admissible populations $q=0$ and $q\equiv\rho$ imply
$\int M_-d\mu\le M$, $\int M_+d\mu\le M/\rho$, and finite
integrals of both work functions. All conditional quantities used below
are thus finite outside a common null set. Define signed conditional
means there, and give them arbitrary values on that null set. For
$f(d)=\mathbb E_{d,+}Y-\mathbb E_{d,-}Y$, the two bias bounds imply
\[
 \int f\,d\mu=(\mathbb E_\rho Y-\mathbb E_0Y)/\rho\ge31/32.
\]
Conditioning on the latent magnitude and then on its sign at the
constant-prior population gives
\[
 V\ge\operatorname{Var}_\rho(Y)
 \ge\rho(1-\rho)\int f^2\,d\mu
 \ge\rho(1-\rho)(31/32)^2\ge2\rho/3.                   \tag{BC5}
\]
This uses only the bias guarantee, not a pointwise unbiased response.
It also shows
\[
                        M\le\tfrac32 V.               \tag{BC6}
\]

\paragraph{Normalize the first observation.}
The event $A=\{T\ge1\}$ is decided by the independent seed before any
observation, so its probability $p$ is independent of the latent state.
It is positive: otherwise the expectations at target values zero and
$\rho$ could not both satisfy the bias bound. The finite constant
$c=\mathbb E[Y\mathbf1_{A^c}]$ is also population independent. Under
the seed law conditioned on $A$, form the active procedure
\[
                  Z=pY+c,\qquad T_Z=T.
\]
It preserves every expectation of $Y$ and always takes at least one
observation. The conditioning is implementable by resampling independent
seeds before observing the coin; its constants depend on the specified
policy and known laws, not on $q,D,S$.

For $p<1$, expansion, Cauchy--Schwarz on $A^c$, and completing the square
give
\begin{align*}
 \mathbb E_qZ^2
 &=p\mathbb E_q[Y^2\mathbf1_A]+2c\mathbb E_qY-c^2\\
 &\le pM+2c\mathbb E_qY-c^2/(1-p)\\
 &\le pM+(1-p)(\mathbb E_qY)^2
 \le m:=pM+2\rho^2.
\end{align*}
For $p=1$, take $c=0$ and the final bound holds directly. Apply (BC5) to
$Z$ to obtain $m\ge2\rho/3$ and
$pM\ge2\rho/3-2\rho^2\ge\rho/6$. Set
$P=MW_0$ and $W_{0Z}=\mathbb E_0T_Z=W_0/p\ge1$. Then
\[
 m\le4pM\le4P,\qquad mW_{0Z}\le4P,
 \qquad P\ge\rho/6.                                   \tag{BC7}
\]
Indeed $2\rho^2\le\rho/2\le3pM$. Every work quantity in (BC7) is
all-negative work.

\paragraph{Select good gaps within the capped global class.}
Now use conditional means and moments for $Z$, keeping the same notation.
The bad set $B_-:=\{M_->1/64\}$ has mass at most $64m$. The set
\[
                    B_+:=\{M_+-M_->2m/\rho\}
\]
has mass less than $2\rho$. Otherwise nonatomicity supplies
$E\subset B_+$ of mass $2\rho$, and the admissible
$q=\tfrac12\mathbf1_E$ gives
\[
 \mathbb E_qZ^2
 =\int M_-d\mu+\tfrac12\int_E(M_+-M_-)d\mu>2m,
\]
contradicting its bound. The response bad set
$B_f:=\{f<3/4\}$ also has mass less than $2\rho$. If not, choose such
an $E$ of mass $2\rho$. Its population has target $\rho$, but
\[
 \mathbb E_{\mathbf1_E/2}Z-\mathbb E_0Z
       =\tfrac12\int_E f\,d\mu<3\rho/4,
\]
whereas the bias bounds force at least $31\rho/32$.

Remove these bad sets and the conditional integrability null sets. The
remaining good set $G$ has
$\mu(G^c)\le64m+4\rho\le70m$. On it,
\[
 |\mathbb E_-Z|\le1/8,\qquad \mathbb E_+Z\ge5/8,
 \qquad M_-\le1/64,\qquad M_+\le1/64+2m/\rho.           \tag{BC8}
\]
The use of sets of mass $2\rho$ is necessary for the cap $q\le1/2$;
the argument does not substitute a pointwise bound $q\le\rho$.

\paragraph{Build a common sign test.}
Let $n=\lceil512(1+m/\rho)\rceil\le1282m/\rho$, the inequality using
$m/\rho\ge2/3$. On the same fixed-gap coin, restart the active procedure
$n$ times with independent seeds and consecutive unused flip blocks.
Output plus if their mean exceeds $3/8$. At a fixed latent state the
restarts are iid: condition on each finite stopping prefix, use the
independence of its unused iid suffix, and sum over the countably many
prefixes. Finite conditional expected work ensures almost-sure
termination for both signs at each good gap. No new latent magnitude is
drawn. The same test is used at every good gap. By (BC8), either mean is
at least $1/4$ from the threshold, and Chebyshev bounds either error by
\[
                 \frac{16(1/64+2m/\rho)}n<1/10.
\]
Its negative expected work at gap $d$ is exactly $nw_-(d)$.

\paragraph{Average negative work over separated annuli.}
For $j\ge3$, put
\[
 I_j=(a4^{-j}/2,a4^{-j}],\qquad
 \mu(I_j)=\tfrac34 16^{-j},\qquad
 J=\left\lfloor\log_{16}\frac3{560m}\right\rfloor,
 \qquad K=J-2.
\]
When $J\ge3$, every retained annulus has mass at least $140m$, so
$G$ occupies at least half of it. Choose $d_j\in G\cap I_j$ with
\[
 w_-(d_j)\le\frac4{\mu(I_j)}\int_{I_j}w_-\,d\mu.
\]
This follows by averaging; the integral is positive since the active
procedure always observes at least once. The choices are existential
proof objects, not inputs to the policy. Successive gap ratios exceed
two, every gap is at most $1/256$, and
\[
 \sum_{j=3}^Jd_j^2w_-(d_j)
       \le\frac{16a^2}{3}\int w_-d\mu
       =\frac{16a^2}{3}W_{0Z}.
\]
Lemma~\ref{lem:finite-sign-entropy}, the restart bound, and (BC7) imply,
for $K\ge K_s$,
\[
 c_sK\log K
 \le\sum_{j=3}^Jd_j^2nw_-(d_j)
 \le\frac{16a^2}{3}\frac{1282m}{\rho}W_{0Z}
 \le\frac{30000a^2P}{\rho}.                            \tag{BC9}
\]
The last numerical constant is conservative, since
$(16/3)\cdot1282\cdot4<30000$.

\paragraph{Obtain one positive constant for every rarity.}
For sufficiently small $\rho$, if $P\ge\sqrt\rho$ the desired order
already follows because $\sqrt\rho LH\to0$. Otherwise (BC7) gives
$m\le4\sqrt\rho$ and
\[
 K\ge\frac{\tfrac12\log(1/\rho)-\log(2240/3)}{\log16}-3.
\]
For all sufficiently small universal $\rho$ this exceeds $K_s$, is at
least a positive universal multiple of $L$, and satisfies
$\log K\ge\tfrac12\log L$. Equation (BC9), with the fixed $a=1/4$,
gives $P\ge c\rho LH$. On the remaining compact interval of rarities,
the lower bound $P\ge\rho/6$ in (BC7) allows the same conclusion after
decreasing one positive universal constant. Finally (BC6) transfers the
raw-moment lower bound to $VW_0$, proving (BC1).
\end{proof}

The all-negative cost is the strengthened part of this statement. It
does not interchange separate suprema, assert a same-population product,
or by itself reduce an arbitrary population algorithm to a single coin.

% END LOCKED SOURCE: base_cost.tex
% BEGIN LOCKED SOURCE: affine_comparison.tex
\section{An RMS comparison for finite convex experiments}
Let $X$ be a finite-dimensional nonempty compact convex set and $\ell(x)$
a linear functional. Consider either of the following experiments:
\begin{enumerate}
\item independent finite categorical observations, whose positive probability
vectors $p_i(x)$ are affine functions of $x$;
\item a finite vector of independent Poisson counts, whose strictly positive
intensities $v_j(x)$ are affine functions of $x$.
\end{enumerate}
In the first case let $\mathcal F$ be the finite-dimensional space of
functions $c+\sum_i\phi_i(z_i)$. In the second let
$\mathcal F=\{c+\sum_j\beta_j z_j\}$. Let $P_x$ denote the whole
observation law. All exponential moments of functions in $\mathcal F$
are finite. For fixed $\phi\in\mathcal F$, the function
$x\mapsto\log\E_xe^\phi$ is concave: it is a sum of logarithms of positive
affine forms in the categorical case, and an affine function of the
intensities in the Poisson case. Moreover $\log(dP_x/dP_y)\in\mathcal F$.

\begin{lemma}[RMS affine comparison]\label{lem:comparison}
For any measurable estimator $G$ with
$r_G=\sup_{x\in X}(\E_x(G-\ell(x))^2)^{1/2}<\infty$ and any $\eta>0$,
there exists $F\in\mathcal F$ such that
\[
 \sup_{x\in X}(\E_x(F-\ell(x))^2)^{1/2}\le12r_G+\eta.       \tag{2}
\]
The result permits nonattainment of the infimum and imposes no restriction
on the original estimator's output range.
\end{lemma}
\begin{proof}
Put $s=\log2$. For $\alpha>0$ and $\phi\in\mathcal F$ define
\begin{align*}
 \Phi(x,y;\phi,\alpha)
 &=\ell(x)-\ell(y)
   +\alpha\log\E_y e^{\phi/\alpha}
   +\alpha\log\E_x e^{-\phi/\alpha}+2\alpha s,\\
 2S&=\inf_{\phi,\alpha>0}\max_{x,y\in X}\Phi(x,y;\phi,\alpha).
\end{align*}
This real-valued continuous function is convex in $(\phi,\alpha)$ by the
perspective transformation and concave in $(x,y)$ by the preceding
log-moment observation. Sion's minimax theorem applies because $X\times X$
is compact convex and the other domain is convex. Thus the infimum and
supremum commute. For fixed $x,y$, Cauchy--Schwarz gives
\[
 \inf_{\psi\in\mathcal F}
  \left(\log\E_y e^\psi+\log\E_x e^{-\psi}\right)
   =2\log\Aff(P_x,P_y),
\]
with equality at $\psi=\tfrac12\log(dP_x/dP_y)$.
Taking the infimum over $\alpha>0$ after putting $\phi=\alpha\psi$
therefore proves
\[
 2S=\max\{\ell(x)-\ell(y):x,y\in X,\ \Aff(P_x,P_y)\ge1/2\}. \tag{3}
\]
The constraint is closed and nonempty, so the maximum exists. Continuity
of affinity is immediate from the finite categorical formula and, for
Poisson vectors, from
$\Aff(P_x,P_y)=\exp[-\tfrac12\sum_j(\sqrt{v_j(x)}-\sqrt{v_j(y)})^2]$.
In particular $0\le S<\infty$.

For the maximizing pair in (3), total variation is at most $\sqrt3/2$
by the standard Cauchy--Schwarz affinity bound
$\|P-Q\|_{\rm TV}\le\sqrt{1-\Aff(P,Q)^2}$.
Testing the two endpoints by thresholding $G$ at their functional midpoint
gives, with $\Delta=2S$,
\[
 2r_G^2\ge \frac{\Delta^2}{4}
           (1-\|P_x-P_y\|_{\rm TV}),\qquad
 S\le\sqrt{\frac{2}{1-\sqrt3/2}}\,r_G.                    \tag{4}
\]
Indeed each wrong decision has squared estimation error at least
$\Delta^2/4$, and the sum of the two testing errors is at least one minus
total variation. The case $S=0$ requires no division by $S$.

Fix $R>S$. Choose $(\phi,\alpha)$ with
$\max_{x,y}\Phi\le2R$. On the diagonal, Cauchy--Schwarz gives
$\Phi(x,x;\phi,\alpha)\ge2\alpha s$, hence $\alpha\le R/s$.
Set
\[
 U=\max_x\{\ell(x)+\alpha\log\E_xe^{-\phi/\alpha}\},\qquad
 V=\max_x\{-\ell(x)+\alpha\log\E_xe^{\phi/\alpha}\}.
\]
Then $U+V+2\alpha s\le2R$. The additive estimator
$F=\phi+(U-V)/2$ satisfies, uniformly in $x$,
\[
 \E_x e^{(F-\ell(x))/\alpha}\le e^{R/\alpha-s},\qquad
 \E_x e^{(\ell(x)-F)/\alpha}\le e^{R/\alpha-s}.
\]
For $t\ge0$, Markov's inequality consequently gives
$\Pp_x(|F-\ell(x)|>R+t)\le e^{-t/\alpha}$, since $2e^{-s}=1$.
Integrating this tail, rather than substituting a confidence guarantee
for an RMS bound, yields
\[
 \E_x(F-\ell(x))^2
 \le R^2+2R\alpha+2\alpha^2
 \le R^2\left(1+\frac2s+\frac2{s^2}\right).                \tag{5}
\]
The product of the squared constants in (4)--(5) is less than $144$:
$\log2>2/3$ and $\sqrt3/2<7/8$ bound it by $8.5\cdot16=136$.
Letting $R\downarrow S$, with the slack chosen for the given $\eta$,
proves (2). This also proves uniform square integrability of the produced
estimator, rather than assuming it.
\end{proof}

\section{From a base-cost obstruction to one pair hard for every capped policy}
In this section only, restrict the parameter set to
\[
 \mathcal X_\rho=\{q:0\le q\le1/2,\ \int q\,d\mu\le\rho\},
 \qquad J=\rho L H.
\]
It is a weak-star compact convex subset of $L^\infty(\mu)$.
The following implication separates its one required entropy input from
the convex-experiment argument. The input is the base-cost bound proved in
the accompanying base-cost section of the full-model note:
\begin{equation}\label{eq:base-cost-input}
 \begin{gathered}
 \text{There is a universal }c_b>0\text{ such that every observable
 single-coin estimator }Y\text{ with}\\
 \sup_{q\in\mathcal X_\rho}|\E_qY-\theta(q)|\le\rho/64,
 \quad V=\sup_{q\in\mathcal X_\rho}\Var_qY<\infty,
 \quad \sup_{q\in\mathcal X_\rho}\E_qT<\infty
 \quad\text{satisfies}\quad V\E_0T\ge c_bJ.
 \end{gathered}                                                     \tag{BC}
\end{equation}
The cost in (BC) is at the all-negative population $q=0$, not the supremum
over populations. The two lemmas below are conditional on precisely this
input. No algorithm-level lower bound is inferred from a variational
optimization alone. Write
$h(P,Q)=1-\Aff(P,Q)$ for squared Hellinger distance in the normalization
used here; it takes values in $[0,1]$.

\subsection*{Exact finite observation reduction}
A capped local policy $A$ includes its independent seed and stops after at
most $m<\infty$ flips of one genuine coin. Its full observed transcript
contains the seed and the stopped word. Its seed distribution and policy
may depend on known parameters, but not on the actual $q$.
For each word $x\in\{0,1\}^t$, let $B_{t,x}$ be the measurable seed set
that causes stopping exactly at that word. For $0\le k\le t\le m$ put
\[
 a_{t,k}=\sum_{x:\,|x|=t,\,\#1(x)=k}\Pp(\text{seed}\in B_{t,x}),
\]
and, with $p_s(d)=1/2+(2s-1)d$, put
\[
 m_{t,k}(q)=\int\big[(1-q(d))p_0(d)^k(1-p_0(d))^{t-k}
                  +q(d)p_1(d)^k(1-p_1(d))^{t-k}\big]\,d\mu.
\]
The finite statistic $(T,K)$, where $K$ is the number of observed heads,
has probabilities $a_{t,k}m_{t,k}(q)$. Conditional on $(T,K)$, the stopped
word and seed have a distribution independent of $q$: their common
likelihood multiplier is exactly $m_{t,k}(q)$ and cancels. This statement
also follows directly by integrating any bounded transcript function over
the finitely many word/seed cells, so it does not assume a latent oracle.
Consequently the full transcript and this finite statistic have exactly
the same Hellinger affinities for all pairs of populations.

Remove coordinates with $a_{t,k}=0$. The remaining coordinates are positive
for every $q$, since $p_s(d)\in[1/4,3/4]$. They are affine weak-star
continuous functions of $q$. Thus their probability vector together with
$\theta(q)$ has a compact convex finite-dimensional range. In particular,
for each capped $A$, its full-transcript squared Hellinger distance
$h_A(q_-,q_+)$ is a continuous convex function of the pair on
$\mathcal X_\rho^2$. Convexity follows from joint concavity of affinity for
two affine probability vectors. This reduction retains the stopping time
and expected cost, not just a fixed number of sampled bits.

\begin{lemma}[A hard pair for each capped policy]\label{lem:hard-pair}
Assume (BC). There are universal constants $c_\delta>0$ and $C_H<\infty$
such that, for every $\rho\in(0,1/4]$, every
$0<\delta\le c_\delta\rho$, and every randomized capped local policy $A$,
there exist $q_-,q_+\in\mathcal X_\rho$ with
$\theta(q_+)-\theta(q_-)\ge\delta$ and
\[
 h_A(q_-,q_+)\le C_H\frac{\delta^2}{J}\E_0T_A.              \tag{6}
\]
The full observed seed and stopped history are included in the Hellinger
distance.
\end{lemma}
\begin{proof}
Let $w=\E_0T_A$. If $w=0$, the seed-only event of taking a first observation
has probability zero, so $A$ has the same transcript law for every $q$.
The pair $q_-=0$, $q_+=\delta$ gives separation $\delta$ and $h_A=0$.
We therefore assume $w>0$ in all the divisions below.
First establish the following additive lower bound for
$n\ge1$ independent runs of $A$. If an additive output with bounded
component functions
$F=c+\sum_{i=1}^nY_i$ has worst RMS $r<\rho/64$, put
$b_i=\E_0Y_i$, $b=\E_0F$, and choose $I$ uniformly in $\{1,\ldots,n\}$
independently. Run just its component on one coin and return
$Z=n(Y_I-b_I)+b$. This preserves every population mean and uses base cost
$w$.

For $\theta=\theta(q)$, the laws of each component output under $q$ and
$0$ have a common subprobability of mass $1-\theta$: integrate the
negative-sign kernel with weight $1-q(d)$. For any two distributions
sharing such a subprobability, writing each as a mixture of its normalized
common part and residual part shows
\[
 (\E_qY_i-\E_0Y_i)^2
 \le\frac{2\theta}{1-\theta}
       (\Var_qY_i+\Var_0Y_i).                              \tag{7}
\]
Indeed the squared displacement of either mixture mean from the common
mean is at most $\theta/(1-\theta)$ times its variance; use
$(u-v)^2\le2u^2+2v^2$. The case $\theta=0$ is identical laws.
Independence and $\sup_q\Var_qF\le r^2$ imply
\[
 \sum_i(\E_qY_i-b_i)^2\le\frac{4\rho}{1-\rho}r^2,
 \quad
 \Var_qZ\le n\sum_i\Var_qY_i
                 +n\sum_i(\E_qY_i-b_i)^2
              \le\frac73nr^2.
\]
All outputs here are bounded, and stopping is capped, so (BC) applies to
$Z$ without an additional integrability assumption. It gives
$r^2\ge3c_bJ/(7nw)$. Including the branch $r\ge\rho/64$ establishes
\[
 r\ge c_a\min\{\rho,\sqrt{J/(nw)}\},\qquad
 c_a=\min\{1/64,\sqrt{3c_b/7}\}.                           \tag{8}
\]

For $w>0$, let $S_n$ be the saddle quantity $S$ from the proof of
Lemma~\ref{lem:comparison} for $n$ independent copies of the finite
statistic of $A$, on the parameter set $\mathcal X_\rho$.
Equations (3) and (5) of that proof show
\[
 2S_n=\max\{\theta(q_+)-\theta(q_-):
             \Aff(P_{A,q_+},P_{A,q_-})^n\ge1/2\},
 \quad \inf_{F\text{ additive}}\sup_q\|F-\theta(q)\|_2\le3S_n.
\]
The latter constant is valid because
$1+2/\log2+2/(\log2)^2<9$. The finite statistic outputs are bounded, so
(8) applies. Hence, for $\beta=2c_a/3$,
\[
 2S_n\ge\beta\min\{\rho,\sqrt{J/(nw)}\}.                 \tag{9}
\]
Take $c_\delta=\beta/2$ and $C_H=8/\beta^2$.
For $0<\delta\le c_\delta\rho$ set
$x=\beta^2J/(4\delta^2w)$. If $x\ge2$, choose $n=\lfloor x\rfloor$,
so $x/2\le n\le x$. Equation (9) gives a pair of separation at least
$2\delta$, with
\[
 h_A\le1-e^{-(\log2)/n}
       \le\frac{\log2}{n}
       \le\frac{8\log2}{\beta^2}\frac{\delta^2w}{J}
       \le C_H\frac{\delta^2w}{J}.
\]
If $x<2$, the pair $q_-=0$, $q_+=\delta$ belongs to
$\mathcal X_\rho$, has separation $\delta$, and $h_A\le1$ suffices
because $C_H\delta^2w/J>1$.
\end{proof}

\begin{lemma}[One pair hard for every capped policy]\label{lem:universal-pair}
Assume (BC), and keep the universal constants of
Lemma~\ref{lem:hard-pair}. For every $0<\rho\le1/4$ and
$0<\delta\le c_\delta\rho$, there exist
$q_-,q_+\in\mathcal X_\rho$, with
$\theta(q_+)-\theta(q_-)\ge\delta$, such that simultaneously for every
randomized capped single-coin policy $A$,
\[
 h_A(q_-,q_+)\le C_H\frac{\delta^2}{\rho LH}\E_0T_A.       \tag{10}
\]
Each policy's seed and full stopped observable history are included.
The constants and pair are chosen before the policy.
\end{lemma}
\begin{proof}
Let $\mathcal K_\delta=\{(q_-,q_+)\in\mathcal X_\rho^2:
\theta(q_+)-\theta(q_-)\ge\delta\}$. It is nonempty compact convex in
the weak-star topology. Fix any finite list of capped policies
$A_1,\ldots,A_k$ and put $w_i=\E_0T_{A_i}$ and
$\lambda=C_H\delta^2/(\rho LH)$. For a simplex vector
$\pi\in\Delta_k$, define the observable mixture policy that draws
$i\sim\pi$, reveals this index in its independent seed, and runs $A_i$.
This is still capped, by the largest of the finitely many caps. Its
base cost is $\sum_i\pi_iw_i$ and, since the revealed labels occupy
disjoint transcript components, its squared Hellinger distance is exactly
$\sum_i\pi_i h_{A_i}$. Lemma~\ref{lem:hard-pair} therefore gives
\[
 \min_{z\in\mathcal K_\delta}
       \sum_i\pi_i\big(h_{A_i}(z)-\lambda w_i\big)\le0
       \qquad(\pi\in\Delta_k).
\]
For the finite simplex the displayed function is continuous affine in
$\pi$ and continuous convex in $z$, by the finite observation reduction.
Sion's theorem now gives
\[
 \min_{z\in\mathcal K_\delta}
       \max_{1\le i\le k}\big(h_{A_i}(z)-\lambda w_i\big)
   =\max_{\pi\in\Delta_k}\min_{z\in\mathcal K_\delta}
       \sum_i\pi_i\big(h_{A_i}(z)-\lambda w_i\big)\le0.
\]
Thus the closed subsets
$\{z\in\mathcal K_\delta:h_A(z)\le\lambda\E_0T_A\}$
have the finite intersection property as $A$ ranges over all capped
policies. Compactness of $\mathcal K_\delta$ gives a point in their full
intersection, proving (10). This uses only finite revealed mixtures and
a finite simplex minimax theorem; no topology on a space of all policies,
unjustified exchange of arbitrary infima, or unobserved-mixture affinity
identity is needed.
\end{proof}


% END LOCKED SOURCE: affine_comparison.tex

\section{From one universal pair to every population algorithm}
\begin{lemma}[Finite-prefix midpoint decomposition]\label{lem:population-prefix}
Suppose $q_-,q_+\in\Xrho$ have this property: every observable randomized
single-coin policy with a deterministic finite cap obeys
\[
 \Hsq(P^A_{q_-},P^A_{q_+})\le\lambda\E_0 T_A,           \tag{B}
\]
where the transcript includes its independent seed and stopped observations,
and $\Hsq(P,Q)=1-\int\sqrt{dP\,dQ}$. Let $\bar q=(q_-+q_+)/2$.
Then every population algorithm capped after $K$ total flips obeys
\[
 \Hsq(P_{q_-}^{\rm pop},P_{q_+}^{\rm pop})
       \le2\lambda\E_{\bar q}T.                       \tag{C}
\]
\end{lemma}
\begin{proof}
Condition first on the original independent seed. Pad the list of observed
coins to $K$ slots by their first-observation indices; an unobserved slot
contributes the empty word of probability one. The policy is a finite binary
decision tree, whose choices are determined by past outcomes and the fixed
seed. For a leaf $z$, let $h_i(z)$ be the ordered word observed from slot $i$.
Under a single fresh coin from $q_r$, write its word probability as
$p_r(h)$, $r\in\{-,+\}$. Independence of the coin tapes gives leaf density
\[
 P_r(z)=\prod_{i=1}^Kp_r(h_i(z)).
\]
Interleaving changes the tree and the words, but not this factorization.
Set $m_i=(p_-(h_i)+p_+(h_i))/2$ and
$c_i=\sqrt{p_-(h_i)p_+(h_i)}$. Word probabilities are affine in $q$, so
$m_i=p_{\bar q}(h_i)$. These quantities are positive; empty words cause
no exception. Since $0\le c_i\le m_i$, the elementary product inequality
gives
\[
 \prod_i m_i-\prod_i c_i
   \le\sum_i(m_i-c_i)\prod_{j\ne i}m_j.                \tag{D}
\]
The left side sums over leaves to $\Hsq(P_-^{\rm pop},P_+^{\rm pop})$,
because the midpoint experiment is normalized. The $i$th right side sums
to the squared Hellinger distance between the two hybrid population
experiments having slot $i$ drawn from $q_-$ or $q_+$ and all other slots
independently drawn from $\bar q$. Those finite experiments are normalized
too. Integrating the original common seed preserves this inequality: joint
Hellinger distance with a revealed common seed is the integral of conditional
Hellinger distances.

For each slot $i$, realize the hybrid experiment by a single-coin policy $A_i$.
Simulate all other slots using independent artificial tapes from the fixed
known reference population $\bar q$, and supply slot $i$ by one genuine
coin. If that slot is never requested, take no genuine observations. Include
the original seed and all artificial tapes in the independent auxiliary seed.
This is an observable single-coin policy capped at $K$ genuine flips.
It may use the fixed reference pair as design constants, but these do not
depend on the genuine coin's unknown endpoint law. The universal quantifier
in (B) covers this policy. No artificial tape is supplied to the original
population algorithm.

The hybrid population transcript is a measurable function of this larger
seed and the stopped genuine observations. Hellinger data processing and
(B) therefore bound its distance by $\lambda\E_0T_{A_i}$.
Since $0\le\bar q\le1/2$, the one-coin latent distribution under $\bar q$
dominates one half of the all-negative distribution, with the magnitude
law fixed. The same domination holds after adjoining the common auxiliary
seed and taking any measurable nonnegative work functional. Hence
\[
 \E_0 T_{A_i}\le2\E_{\bar q}T_{A_i}.
\]
When the genuine coin too is from $\bar q$, its expected charged work is
exactly the work on slot $i$ in the original midpoint experiment. Summing
over the at most $K$ slots yields (C). All simulations are only devices in
the lower-bound comparison; every genuine flip of the population algorithm
is counted exactly once.
\end{proof}

\begin{proof}[Proof of the lower bound in Theorem~\ref{thm:critical-minimax}]
Take the universal constants $c_\delta,C_H$ from
Lemma~\ref{lem:universal-pair}. Choose a fixed
\[
 0<\gamma\le\min\{c_\delta,1/(4\sqrt{C_H})\},
 \qquad \delta=\gamma\min\{\rho,\sqrt{J/B}\}.
\]
That lemma provides one pair $q_-,q_+\in\Xrho$ with target separation at
least $\delta$ and property (B) with $\lambda=C_H\delta^2/J$.
The pair is valid simultaneously for all finite caps and policies; it is
not reselected after seeing a population transcript.

Let $A$ be any full population algorithm with uniform expected work at most
$B$. Apply Lemma~\ref{lem:population-prefix} to its prefix capped after
$K$ flips, recording whether it already halted and, when relevant, its
terminal output as a measurable statistic. Since $\bar q\in\Xrho\subset
\mathcal Q_\rho$, its midpoint expected capped work is at most $B$.
Thus every prefix has squared Hellinger distance at most
$2C_H\delta^2B/J\le1/8$. The inequality
$\operatorname{TV}(P,Q)\le\sqrt{2\Hsq(P,Q)}$ bounds its total variation
by $1/2$.

Uniform expected work implies almost-sure termination separately under both
endpoints. For any terminal event $E$, the event $E\cap\{T\le K\}$ is
determined by the capped prefix, and the omitted endpoint probabilities
are at most $\Pp_{q_-}(T>K)+\Pp_{q_+}(T>K)\to0$.
Consequently the terminal transcript's total variation is also at most
$1/2$. This passage assumes neither a globally bounded number of coins
nor differentiability or a chain rule at an unbounded stopping time.

If the estimator has infinite worst-case squared error, the lower bound is
automatic. Otherwise threshold its output halfway between the endpoint
targets. With endpoint separation $\Delta\ge\delta$ and worst RMS $r$,
Markov's inequality bounds the sum of the two testing errors by
$8r^2/\Delta^2$. Every test has error sum at least
$1-\operatorname{TV}\ge1/2$. Hence $r\ge\Delta/4\ge\delta/4$,
which proves the lower inequality in (A), including saturation.
\end{proof}

\section{Consequence for the proposed rare-event modulus}
For this exact law the variational optimizer gives a cutoff $t\in(0,a)$:
\[
 \rho=(t/a)^2[1+2\log(a/t)],\qquad
 \Psi_{\mu,\rho}(h)=
 \begin{cases}
 [2ht(1+\log(a/t))-h^2]/a^2,&0<h<t,\\
 \rho,&h\ge t.
 \end{cases}                                          \tag{E}
\]
These are inherited variational identities, not an attaining statistical
algorithm. Put $u=1+2\log(a/t)$. Then $\rho=u e^{1-u}$ and
$L=u-\log u$, so $L\le u\le2L$. For $h<t$,
$\Psi(h)=(ht/a^2)(u+1-h/t)$; its coefficient lies between $u$ and $u+1$.
Together with the saturated branch this proves, with universal constants
and fixed $a=1/4$,
\[
 \Psi_{\mu,\rho}(B^{-1/2})\asymp\min\{\rho,\sqrt{\rho L/B}\}.
\]
For instance take $\rho_n=\exp(-\exp n)$ and $B_n=\rho_n^{-2}$.
Both radii are eventually unsaturated, and (A) implies
\[
 \frac{R_{\mu,\rho_n}(B_n)}{\Psi_{\mu,\rho_n}(B_n^{-1/2})}
       \ge c'\sqrt{1+\log(1+\log(1/\rho_n))}\longrightarrow\infty.
\]
Thus no one finite universal $C$ can
satisfy the conjectured upper bound even at this single known law, with
the algorithm knowing $\rho$ and $B$. The missing statistical cost is not
hidden magnitudes alone: known-profile submodels can behave differently.
The lower bound uses the entire convex nuisance class to obtain a common
hard pair. It does not assert that the radial variational optimizer itself
is that pair, nor give a closed-form expression for the pair selected by
compactness. No claim for every other magnitude law follows from (A).

\appendix
% BEGIN LOCKED SOURCE: upper_bound.tex
\section{Self-contained upper checkpoint and provenance}
This appendix records the existing upper argument needed for the upper
inequality in Theorem~\ref{thm:critical-minimax}, so the characterization does not depend on an inaccessible
working-tree file. The inspected source was
\nolinkurl{research/2026-10-27-minimax-audit/theorem/critical.tex}, SHA-256
\begin{center}\small\nolinkurl{0e8b9eaf762b3b09fa0d92ce1a036c2ceea8c12d00a5eea8286b852f36a81536}\end{center}
It was read without modifying it. The following is a compact restatement of
that pre-existing proof, not a new construction or discovery claim.
The source is now committed at
\nolinkurl{a8a8dc49115aabdd1fba8d53d3390eb0a5eee864}.

Put $H=1+\log L$. For $j\ge0$ choose
\[
 \varepsilon_j=16^{-(j+1)},\qquad
 N_j=\lceil(a^2\rho\varepsilon_j)^{-1}\rceil,\qquad
 K_j=\lceil\log_2(a\sqrt{N_j})\rceil.
\]
For $0\le k\le K_j$, let $r_k=a2^{-(k+1)}$ and, after $n$ observations
with $h_n$ heads, define
\[
 L^+_{k,n}=(1+2r_k)^{h_n}(1-2r_k)^{n-h_n},\qquad
 L^-_{k,n}=(1-2r_k)^{h_n}(1+2r_k)^{n-h_n}.
\]
Before $N_j$, stop with output zero if a negative likelihood reaches
$(K_j+1)/\varepsilon_j$, or with output one if a positive likelihood reaches
$(K_j+1)/(\rho\varepsilon_j)$. Fix zero priority for ties. At $N_j$, if no
earlier decision was made, return the strict majority bit. Denote this bounded,
observable classifier by $Z_j$ and its work by $T_j$.

For the wrong-sign likelihood, each next multiplier has mean $1-4r_kd\le1$.
Stopping this nonnegative supermartingale at a crossing and a deterministic
horizon, and then increasing the horizon, bounds the probability of crossing
level $b$ by $1/b$. Summing over $K_j+1$ references gives wrong-crossing
probabilities at most $\varepsilon_j$ on a positive coin and
$\rho\varepsilon_j$ on a negative coin. The wrong forced-majority event is
a subset of the full-tape event, of probability at most $e^{-2N_jd^2}$.
This is Hoeffding's bound: the centered Bernoulli log moment-generating
function has second derivative at most $1/4$, so is at most $\lambda^2/8$;
use independence and exponential Markov with $\lambda=\pm4d$.
For the exact power law,
\[
 \E_\mu e^{-2ND^2}=\frac{1-e^{-2Na^2}}{2Na^2}\le\frac1{2Na^2}.
\]
No independence of a stopping event and this full-tape majority event is used.
It follows, uniformly over the global-rarity class, that
\[
 \Pp(Z_j\ne S)\le\tfrac52\rho\varepsilon_j,
 \qquad \Pp(Z_j=1,S=0)\le\tfrac32\rho\varepsilon_j.       \tag{U1}
\]

For $d\ge N_j^{-1/2}$, choose $k=\lceil\log_2(a/d)\rceil\le K_j$.
Then $d/4\le r_k\le d/2$, and the correct-sign log likelihood has mean
increment at least $2d^2/3$: use
\[
 \tfrac12\log(1-4r_k^2)
       +d\log\frac{1+2r_k}{1-2r_k}
 \ge-\tfrac83r_k^2+4dr_k\ge\tfrac23d^2.
\]
Its upward increment is less than one. The bounded stopped-sum identity for
the crossing time $\sigma$ at a finite horizon $n$ gives
$g\E(\sigma\wedge n)\le A_{j,s}$, where
\[
 A_{j,s}=1+\log(K_j+1)+\log(1/\varepsilon_j)+s\log(1/\rho).
\]
Monotone convergence proves finite expected crossing time and, including the
cap and the trivial small-gap case,
\[
 \E(T_j\mid d,s)\le\min\{N_j,(3/2)A_{j,s}/d^2\}.        \tag{U2}
\]
Write $c_0=3/2$. Direct integration of the negative contribution gives
\[
 \E\min\{N,c_0A/D^2\}
 \le(c_0A/a^2)[1+\log_+(Na^2)].
\]
For positives, layer cake and
$\E[q\ind_E]\le\min\{\rho,\mu(E)\}$ give
\[
 \E[q\min\{N,c_0A/D^2\}]
 \le\int_0^\rho\min\{N,c_0A/(a^2u)\}\,du
 \le(c_0A/a^2)[1+\log_+(Na^2\rho)].
\]
The exact integral is $N\rho$ when $Na^2\rho\le c_0A$, and otherwise is
$(c_0A/a^2)[1+\log(Na^2\rho/(c_0A))]$. Thus this step also enforces the
global, rather than pointwise, rarity condition.

Elementary bounds from the definitions are
\[
\begin{gathered}
 N_ja^2\le2/(\rho\varepsilon_j),\qquad K_j+1\le5(L+j+1),\\
 A_{j,0}\le7H(j+1),\qquad A_{j,1}\le7H(j+1)+L,\\
 1+\log(N_ja^2)\le5L(j+1),\qquad
 1+\log(N_ja^2\rho)\le5(j+1).
\end{gathered}
\]
For example, $\log(L+j+1)\le\log L+j+1$ and $4\log2<3$ prove the
bound on $A_{j,0}$. Substitution in (U2) and the two integrals gives
\[
 \E T_j\le(225/2)a^{-2}LH(j+1)^2.                        \tag{U3}
\]

Independently choose a level $J$ with $p_j=2^{-(j+1)}$. For $J=0$ return
$Y=Z_0/p_0$ on one fresh coin. For $J=j\ge1$, run $Z_j$ and $Z_{j-1}$
on successive unused blocks of that same coin, set $\Delta_j=Z_j-Z_{j-1}$,
and return $Y=\Delta_j/p_j$. The bounded stopping times make the iid suffix
argument elementary, as in the restart argument above; both decisions share
the same latent sign. A disagreement implies an error of one classifier.
Hence (U1) gives
\[
 \E Z_0^2\le\tfrac{35}{32}\rho,
 \qquad \E\Delta_j^2\le\tfrac{85}{32}\rho16^{-j}\quad(j\ge1).
\]
Tonelli and the exact geometric sum now give
\[
 \E Y^2\le\frac{35}{16}\rho+
                 \frac{85}{16}\rho\sum_{j\ge1}8^{-j}
            =\frac{165}{56}\rho.
\]
The disagreement bounds also give absolute summability of expected increments,
so signed expectations may be telescoped after integrability is established:
$\E Y=\lim_j\E Z_j=\theta$. Work costs at most twice (U3), and
$\sum_jp_j(j+1)^2=6$, yielding
\[
 \E T_Y\le1350a^{-2}LH.
\]
All operations are on one observable coin with independent seed; the random
level is unbounded but the separate expectations are finite. This completes
the upper checkpoint used in Theorem~\ref{thm:critical-minimax}.
Its actual work also has genuinely unbounded support: at any fixed observation
count $n$, all likelihoods are at most $(1+a)^n$. For sufficiently large $j$,
both the smallest stopping boundary and $N_j$ exceed the relevant finite
thresholds, so $T_j>n$ on every tape. Since $\Pp(J=j)>0$, actual work exceeds
every finite cap with positive probability while retaining finite mean.


% END LOCKED SOURCE: upper_bound.tex

\begin{proof}[Proof of the upper bound in Theorem~\ref{thm:critical-minimax}]
The preceding construction supplies a one-coin unbiased output with
$\sup_q\E_qY^2\le(165/56)\rho$ and
$\sup_q\E_qT_Y\le W_*=1350a^{-2}LH$. Independently select
$N\sim\operatorname{Poisson}(B/W_*)$ and return
$\widehat\theta=(W_*/B)\sum_{i=1}^NY_i$ from independent genuine coins.
Conditioning on $N$ proves unbiasedness, uniform expected work at most $B$,
and variance at most $(165/56)\rho W_*/B$. All moments are finite before
using these identities, by the bounds and the finite Poisson moments.
Choose the zero-work constant $\rho/2$ whenever its known error bound is
smaller. Both choices are observable and depend only on known parameters.
With $a$ fixed, this proves the upper inequality in (A). This established
replication device is not claimed as new.
\end{proof}

\begin{thebibliography}{9}
\bibitem{jn} A. Juditsky and A. Nemirovski.
\emph{Nonparametric estimation by convex programming}.
Annals of Statistics 37 (2009), 2278--2300.
\url{https://www2.isye.gatech.edu/~nemirovs/AnnStat_Saddle_2009.pdf}.
\bibitem{pw} Y. Polyanskiy and Y. Wu.
\emph{Dualizing Le Cam's method for functional estimation, with applications
to estimating the unseens}. arXiv:1902.05616v2 (2021), Theorem 1,
Proposition 4 and Section 4.1.
\url{https://arxiv.org/abs/1902.05616}.
\bibitem{lv} J. C. H. Lee and P. Valiant.
\emph{Uncertainty about Uncertainty: Optimal Adaptive Algorithms for
Estimating Mixtures of Unknown Coins}. SODA 2021; arXiv:1904.09228v3.
\url{https://arxiv.org/html/1904.09228v3}.
\bibitem{clq} L. Chen, J. Li and M. Qiao.
\emph{Towards Instance Optimal Bounds for Best Arm Identification}.
COLT 2017; arXiv:1608.06031v2.
\url{https://arxiv.org/html/1608.06031v2}.
\bibitem{critical} Existing workspace critical upper construction,
\nolinkurl{research/2026-10-27-minimax-audit/theorem/critical.tex},
commit \nolinkurl{a8a8dc49115aabdd1fba8d53d3390eb0a5eee864}.
\end{thebibliography}
\end{document}
